\documentclass[11pt]{article}
\usepackage[a4paper,margin=1in]{geometry}
\usepackage{amsmath,amssymb,amsthm}
\usepackage{hyperref}

\newcommand{\dd}{\mathrm{d}}
\newcommand{\E}{\mathbb{E}}
\newcommand{\eps}{\varepsilon}
\newcommand{\cL}{\mathcal{L}}
\newcommand{\cT}{\mathcal{T}}
\renewcommand{\Re}{\operatorname{Re}}
\newcommand{\nrm}[2]{\|#1\|_{#2}}

\newtheorem{theorem}{Theorem}
\newtheorem{lemma}[theorem]{Lemma}
\newtheorem{proposition}[theorem]{Proposition}
\newtheorem{corollary}[theorem]{Corollary}
\theoremstyle{remark}
\newtheorem{remark}[theorem]{Remark}

\title{Orbit decorrelation at $N^{1/3}$:\\
a proved model case of the plasma closure strategy}
\author{Working note}
\date{16 November 2026}

\begin{document}
\maketitle

\section{What is and is not proved}

The strategy in \texttt{plasma\_extension.tex} has eight steps. This note
proves two things and no more.
\begin{itemize}
  \item \textbf{Step 1 in full} (Lemma~\ref{lem:fluct}): the fluctuation
    system is exact for the $N$-particle dynamics.
  \item \textbf{Step 5 in a model} (Theorem~\ref{thm:main}): for a tagged
    particle driven by a \emph{prescribed} stationary Gaussian field of
    amplitude $N^{-1/2}$ with Markovian mode amplitudes, the departure from
    the ballistic orbit has characteristic function
    $\exp(-t^3\,q\cdot Dq/(3N))$ up to an explicit error, which is
    $O(N^{-1/6})$ at $t\sim N^{1/3}$. The proof is non-perturbative: it does
    not expand the orbit.
\end{itemize}
The full strategy is \emph{not} proved. In the model the field is given in
advance; in the particle system it is generated by the particles, including
the tagged one, and is neither Gaussian nor Markovian. Steps 4 and 6--8 are
untouched. Section~\ref{sec:gap} lists what separates the model from the real
problem. The method of proof (perturbed test functions for a fast Markov
process) is classical; the result is a model computation, not a new
technique.

\section{The exact fluctuation system}

Let $\phi\in C^2(\mathbb{T}^d)$ be even. Consider
\begin{equation}
  \dot x_i = v_i, \qquad
  \dot v_i = -\frac1N\sum_{j\ne i}\nabla\phi(x_i - x_j),
  \qquad i = 1,\dots,N,
  \label{eq:nbody}
\end{equation}
with random initial data such that $\E|v_i(0)| < \infty$. Let
$\mu_t = N^{-1}\sum_i\delta_{(x_i(t),v_i(t))}$, $f_t = \E\mu_t$,
$\eta_t = \sqrt N(\mu_t - f_t)$, and for a finite signed measure $h$ put
$F[h](x) = -\int\nabla\phi(x - x')\,\dd h(x',v')$.

\begin{lemma}\label{lem:fluct}
For every $\psi\in C_c^\infty(\mathbb{T}^d\times\mathbb{R}^d)$,
\begin{align}
  \frac{\dd}{\dd t}\langle f,\psi\rangle
   &= \langle f,\, v\cdot\nabla_x\psi + F[f]\cdot\nabla_v\psi\rangle
     + \frac1N\,\E\langle\eta,\,F[\eta]\cdot\nabla_v\psi\rangle,
     \label{eq:mean}\\
  \frac{\dd}{\dd t}\langle\eta,\psi\rangle
   &= \langle\eta,\, v\cdot\nabla_x\psi + F[f]\cdot\nabla_v\psi\rangle
     + \langle f,\,F[\eta]\cdot\nabla_v\psi\rangle \notag\\
   &\quad + \frac{1}{\sqrt N}\Bigl(\langle\eta,F[\eta]\cdot\nabla_v\psi\rangle
       - \E\langle\eta,F[\eta]\cdot\nabla_v\psi\rangle\Bigr).
     \label{eq:fluct}
\end{align}
If the initial law is exchangeable, $f$ is the one-particle marginal.
\end{lemma}

\begin{proof}
Since $\phi$ is even and $C^1$, $\nabla\phi(0) = 0$, so the restriction
$j\ne i$ in \eqref{eq:nbody} may be dropped and
$\dot v_i = F[\mu_t](x_i)$. Hence
\[
  \frac{\dd}{\dd t}\langle\mu,\psi\rangle
  = \frac1N\sum_i\bigl(v_i\cdot\nabla_x\psi + F[\mu](x_i)\cdot\nabla_v\psi\bigr)(x_i,v_i)
  = \langle\mu,\,v\cdot\nabla_x\psi + F[\mu]\cdot\nabla_v\psi\rangle .
\]
All terms are integrable because $\nabla\phi$ and the derivatives of $\psi$
are bounded and $\psi$ has compact support. The map
$(\mu,\nu)\mapsto\langle\mu,F[\nu]\cdot\nabla_v\psi\rangle$ is bilinear.
Substituting $\mu = f + N^{-1/2}\eta$ and using $\E\eta = 0$,
\[
  \E\langle\mu,F[\mu]\cdot\nabla_v\psi\rangle
  = \langle f,F[f]\cdot\nabla_v\psi\rangle
    + \frac1N\E\langle\eta,F[\eta]\cdot\nabla_v\psi\rangle ,
\]
which gives \eqref{eq:mean} after taking expectations. Subtracting
\eqref{eq:mean} from the equation for $\mu$ and multiplying by $\sqrt N$ gives
\eqref{eq:fluct}. Under exchangeability
$\E\langle\mu,\psi\rangle = \E\psi(x_1,v_1)$.
\end{proof}

\section{The model}

Fix a finite set $K\subset\mathbb{R}^d\setminus\{0\}$ and, for $k\in K$,
vectors $f_k\in\mathbb{C}^d$ and numbers $\gamma_k>0$, $\omega_k\in\mathbb{R}$.
Let $(B_k)_{k\in K}$ be independent complex Brownian motions with
$\E|B_k(t)|^2 = t$ and $\E B_k(t)^2 = 0$, and let $a_k$ be the stationary
solutions of
\[
  \dd a_k = -(\gamma_k + i\omega_k)\,a_k\,\dd\tau + \sqrt{2\gamma_k}\,\dd B_k ,
\]
so that $\E|a_k|^2 = 1$, $\E a_k^2 = 0$. Define the random field
\begin{equation}
  F(x,\tau) = \sum_{k\in K}\Re\bigl(f_k\,a_k(\tau)\,e^{ik\cdot x}\bigr).
\end{equation}
It is Gaussian, stationary in space and time, with mean zero. The choice
$f_k = -i\varphi_k k$ makes it a potential field. For $\eps>0$ a particle
starting from $(x_0,v_0)$ obeys
\begin{equation}
  \dot x = v, \qquad \dot v = \eps\,F(x,\tau).
\end{equation}
The correspondence with the particle system is $\eps = N^{-1/2}$.

Write $Y = x - x_0 - v_0\tau$ and $w = v - v_0$ for the departure from the
ballistic orbit, and $b_k(\tau) = a_k(\tau)e^{ik\cdot(x_0+v_0\tau)}$. Then
$b$ is again a stationary Ornstein--Uhlenbeck process,
\begin{equation}
  \dd b_k = -\lambda_k b_k\,\dd\tau + \sqrt{2\gamma_k}\,\dd\tilde B_k,
  \qquad \lambda_k = \gamma_k + i(\omega_k - k\cdot v_0),
\end{equation}
with $\tilde B_k$ complex Brownian motions of the same law, and
\begin{equation}
  \dot Y = w, \qquad \dot w = \eps\,G(Y,b), \qquad
  G(Y,b) = \sum_{k\in K}\Re\bigl(f_k b_k e^{ik\cdot Y}\bigr),
  \qquad Y_0 = 0,\ w_0 = 0 .
  \label{eq:Yw}
\end{equation}
The process $(Y,w,b)$ is a diffusion with generator
\begin{equation}
  \cL = w\cdot\nabla_Y + \eps\,G\cdot\nabla_w + A,
  \qquad
  A = \sum_{k\in K}\Bigl(-\lambda_k b_k\partial_{b_k}
      - \bar\lambda_k\bar b_k\partial_{\bar b_k}
      + 2\gamma_k\,\partial_{b_k}\partial_{\bar b_k}\Bigr).
\end{equation}
Two identities are used repeatedly:
\begin{equation}
  A\,b_k = -\lambda_k b_k,\qquad
  A(b_kb_l) = -(\lambda_k+\lambda_l)\,b_kb_l,\qquad
  A(b_k\bar b_l - \delta_{kl}) = -(\lambda_k+\bar\lambda_l)(b_k\bar b_l - \delta_{kl}).
  \label{eq:Aid}
\end{equation}
Define the diffusion matrix
\begin{equation}
  D = \frac12\sum_{k\in K}\Re\!\Bigl(\frac{1}{\lambda_k}\Bigr)\,
      \Re\bigl(f_k\otimes\bar f_k\bigr),
  \qquad
  \Re\!\Bigl(\frac{1}{\lambda_k}\Bigr)
  = \frac{\gamma_k}{\gamma_k^2 + (\omega_k - k\cdot v_0)^2} .
  \label{eq:D}
\end{equation}
$D$ is symmetric and positive semidefinite, and
$D = \int_0^\infty \E\bigl[F(v_0u,u)\otimes F(0,0)\bigr]_{\mathrm{sym}}\dd u$.
It is a sum of resonances $\omega_k = k\cdot v_0$ broadened by $\gamma_k$.

For a function $h(Y,b)$ write
$\nrm{h}{p} = \bigl(\E\sup_Y|h(Y,b)|^p\bigr)^{1/p}$, the expectation being
over the stationary law of $b$. All functions below are polynomials in
$(b,\bar b)$ with coefficients bounded in $Y$, so these norms are finite.

\section{The theorem}

\begin{theorem}[Orbit decorrelation in a prescribed field]\label{thm:main}
For every $q\in\mathbb{R}^d$, $\eps>0$ and $t\ge 0$,
\begin{equation}
  \Bigl|\,\E\,e^{iq\cdot Y_t} - \exp\!\Bigl(-\tfrac13\,\eps^2t^3\,q\cdot Dq\Bigr)\Bigr|
  \;\le\; C_1\,\eps t + C_2\,\eps^2t^2 + C_3\,\eps^3t^4 + C_4\,\eps^4t^5 ,
  \label{eq:main}
\end{equation}
with constants $C_i$ depending only on $q$, $K$, $(f_k)$ and $(\lambda_k)$,
given explicitly in \eqref{eq:consts}. In particular, with $\eps = N^{-1/2}$
and $t = sN^{1/3}$,
\[
  \E\,e^{iq\cdot Y_t} = \exp\!\Bigl(-\tfrac13\,s^3\,q\cdot Dq\Bigr)
  + O\bigl((s + s^4)\,N^{-1/6}\bigr),
\]
and the estimate is non-trivial for $1\ll t\ll N^{3/8}$.
\end{theorem}

\begin{corollary}\label{cor:law}
As $N\to\infty$ with $t = sN^{1/3}$, $Y_t$ converges in law to a centred
Gaussian vector with covariance $\tfrac23 s^3 D$.
\end{corollary}

\begin{proof}
Pointwise convergence of characteristic functions and L\'evy's continuity
theorem.
\end{proof}

\begin{corollary}[Erasure of phase-mixed information in the mean]\label{cor:mem}
For the particle started at $(x_0,v_0)$,
$\E\,e^{-iq\cdot x_t} = e^{-iq\cdot(x_0+v_0t)}
\bigl[\exp(-\tfrac13\eps^2t^3\,q\cdot D(v_0)q) + \text{error}\bigr]$ with the
error of \eqref{eq:main}. The free-streaming phase survives only for
$t\lesssim(N/q\cdot Dq)^{1/3}$.
\end{corollary}

\section{Proof of Theorem~\ref{thm:main}}

Fix $q$ and $t$, and write $\theta = t-\tau$, $d = q\cdot Dq$.

\subsection{The limiting test function}

Let
\begin{equation}
  U(\tau,Y,w) = \exp\!\Bigl(iq\cdot Y + i\theta\,q\cdot w
     - \tfrac13\,\eps^2\theta^3 d\Bigr).
\end{equation}
Then $|U|\le1$, $U(t,Y,w) = e^{iq\cdot Y}$,
$U(0,0,0) = \exp(-\tfrac13\eps^2t^3d)$, and with
$\cT = \partial_\tau + w\cdot\nabla_Y$,
\begin{equation}
  \cT U = \eps^2\theta^2 d\,U, \qquad \nabla_wU = i\theta\,q\,U .
  \label{eq:Uid}
\end{equation}
($U$ solves the backward equation of the Kolmogorov diffusion
$\dd Y = w\,\dd\tau$, $\dd w = \eps\sqrt{2D}\,\dd W$.)

\subsection{Correctors}

Put $g = q\cdot G$ and define, as functions of $(Y,b)$,
\begin{align*}
  \alpha &= \sum_k\Re\Bigl(\frac{q\cdot f_k}{\lambda_k}\,b_ke^{ik\cdot Y}\Bigr),
  &
  \alpha_2 &= \sum_k\Re\Bigl(\frac{q\cdot f_k}{\lambda_k^2}\,b_ke^{ik\cdot Y}\Bigr).
\end{align*}
By \eqref{eq:Aid}, $A\alpha = -g$ and $A\alpha_2 = -\alpha$. Writing
$A_k = q\cdot f_k/\lambda_k$, $B_k = q\cdot f_k$, $e_k = e^{ik\cdot Y}$,
\[
  \alpha g = \frac14\sum_{k,l}\Bigl(A_kB_l\,b_kb_l\,e_ke_l
     + A_k\bar B_l\,b_k\bar b_l\,e_k\bar e_l + \text{c.c.}\Bigr),
\]
whose mean under the stationary law is
$\frac12\sum_k\Re(A_k\bar B_k) = \frac12\sum_k|q\cdot f_k|^2\Re(1/\lambda_k) = d$.
Define
\begin{equation}
  \beta = -\frac14\sum_{k,l}\Bigl(
     \frac{A_kB_l}{\lambda_k+\lambda_l}\,b_kb_l\,e_ke_l
     + \frac{A_k\bar B_l}{\lambda_k+\bar\lambda_l}\,(b_k\bar b_l-\delta_{kl})\,e_k\bar e_l
     + \text{c.c.}\Bigr).
\end{equation}
By \eqref{eq:Aid}, $A\beta = \alpha g - d$, and $\E\beta = 0$ for each $Y$.
The denominators have real part $\gamma_k+\gamma_l>0$. Set
\begin{equation}
  U^\eps = U\Bigl[\,1 + i\eps\theta\,\alpha
     + i\eps\theta\,(w\cdot\nabla_Y\alpha_2) + \eps^2\theta^2\beta\Bigr].
\end{equation}

\subsection{The remainder}

\begin{lemma}\label{lem:rem}
$(\partial_\tau + \cL)U^\eps = U\,R$, where
\begin{align*}
  R ={}& -i\eps\,\alpha + i\eps^3\theta^3 d\,\alpha
     - i\eps\,(w\cdot\nabla\alpha_2)
     + i\eps\theta\,(w\cdot\nabla)^2\alpha_2
     + i\eps^3\theta^3 d\,(w\cdot\nabla\alpha_2) \\
   & + i\eps^2\theta\,(G\cdot\nabla\alpha_2)
     - \eps^2\theta^2\,(w\cdot\nabla\alpha_2)\,g
     - 2\eps^2\theta\,\beta
     + \eps^2\theta^2\,(w\cdot\nabla\beta)
     + \eps^4\theta^4 d\,\beta
     + i\eps^3\theta^3\,\beta g .
\end{align*}
\end{lemma}

\begin{proof}
Use $\partial_\tau+\cL = \cT + \eps G\cdot\nabla_w + A$, the identities
\eqref{eq:Uid}, $\partial_\tau\theta = -1$, and the fact that $A$ acts only on
$b$ and commutes with $\nabla_Y$.

\emph{$U$:} $(\cT + \eps G\cdot\nabla_w)U = \eps^2\theta^2d\,U + i\eps\theta g\,U$.

\emph{$P_1 = i\eps\theta\alpha U$:}
$\cT P_1 = [-i\eps\alpha + i\eps\theta(w\cdot\nabla\alpha)
+ i\eps^3\theta^3d\,\alpha]U$;
$\eps G\cdot\nabla_wP_1 = -\eps^2\theta^2\alpha g\,U$;
$AP_1 = -i\eps\theta g\,U$.

\emph{$P_2 = i\eps\theta(w\cdot\nabla\alpha_2)U$:}
$\cT P_2 = [-i\eps(w\cdot\nabla\alpha_2)
+ i\eps\theta(w\cdot\nabla)^2\alpha_2
+ i\eps^3\theta^3d\,(w\cdot\nabla\alpha_2)]U$;
$\eps G\cdot\nabla_wP_2 = [i\eps^2\theta(G\cdot\nabla\alpha_2)
- \eps^2\theta^2(w\cdot\nabla\alpha_2)g]U$;
$AP_2 = -i\eps\theta(w\cdot\nabla\alpha)U$.

\emph{$P_3 = \eps^2\theta^2\beta U$:}
$\cT P_3 = [-2\eps^2\theta\beta + \eps^2\theta^2(w\cdot\nabla\beta)
+ \eps^4\theta^4d\,\beta]U$;
$\eps G\cdot\nabla_wP_3 = i\eps^3\theta^3\beta g\,U$;
$AP_3 = \eps^2\theta^2(\alpha g - d)U$.

Three cancellations occur on summing: $i\eps\theta g\,U$ against $AP_1$;
$i\eps\theta(w\cdot\nabla\alpha)U$ against $AP_2$; and
$\eps^2\theta^2d\,U - \eps^2\theta^2\alpha g\,U$ against $AP_3$. The
remaining terms are those listed.
\end{proof}

\subsection{A priori bound and Dynkin's formula}

\begin{lemma}\label{lem:apriori}
Let $\Gamma(b) = \sum_k|f_k||b_k|$ and $M_p = (\E\Gamma^p)^{1/p}$. Then
$(\E|w_\tau|^p)^{1/p}\le\eps\tau M_p$ for $p\ge1$.
\end{lemma}

\begin{proof}
$|G(Y,b)|\le\Gamma(b)$, so $|w_\tau|\le\eps\int_0^\tau\Gamma(b_\sigma)\dd\sigma$.
Apply Minkowski's integral inequality and stationarity of $b$.
\end{proof}

\begin{lemma}\label{lem:dynkin}
$\E\,U^\eps(t,Y_t,w_t,b_t) - \E\,U^\eps(0,0,0,b_0)
= \E\int_0^t (UR)(\tau,Y_\tau,w_\tau,b_\tau)\,\dd\tau$.
\end{lemma}

\begin{proof}
The system \eqref{eq:Yw} has a unique global solution for every path of $b$,
since $G$ is smooth and $|G|\le\Gamma(b_\tau)$ is locally bounded in time.
$U^\eps$ is smooth, and It\^o's formula gives
$U^\eps(t) - U^\eps(0) = \int_0^t(\partial_\tau+\cL)U^\eps\,\dd\tau + \mathcal{M}_t$,
with $\mathcal{M}$ a stochastic integral against the $\tilde B_k$ whose
integrands are $\sqrt{2\gamma_k}\,\partial_{b_k}U^\eps$ and conjugates. These
are bounded by polynomials in $|b|$ and $|w|$, hence square integrable
uniformly on $[0,t]$ by Lemma~\ref{lem:apriori} and Gaussianity of $b$. So
$\mathcal{M}$ is a martingale with zero mean.
\end{proof}

\subsection{Conclusion}

At $\tau=t$, $\theta=0$ and $U^\eps = e^{iq\cdot Y_t}$. At $\tau=0$, $w=0$ and
$Y=0$, so $U^\eps = U(0,0,0)[1 + i\eps t\,\alpha(0,b_0) + \eps^2t^2\beta(0,b_0)]$;
since $\E\alpha = \E\beta = 0$,
$\E\,U^\eps(0) = \exp(-\tfrac13\eps^2t^3d)$ exactly. By Lemma~\ref{lem:dynkin}
and $|U|\le1$,
\[
  \Bigl|\E e^{iq\cdot Y_t} - \exp(-\tfrac13\eps^2t^3d)\Bigr|
  \le \int_0^t\E|R|\,\dd\tau .
\]
Each term of $R$ is bounded using Lemma~\ref{lem:apriori}, the
Cauchy--Schwarz inequality, stationarity of $b$, and
$\int_0^t\theta^a\tau^c\,\dd\tau = t^{a+c+1}\,a!\,c!/(a+c+1)!$:
\begin{center}
\renewcommand{\arraystretch}{1.3}
\begin{tabular}{ll}
  term of $R$ & bound on its time integral \\ \hline
  $\eps\alpha$ & $\eps t\,\nrm{\alpha}{1}$ \\
  $\eps^3\theta^3d\,\alpha$ & $\tfrac14\eps^3t^4\,d\,\nrm{\alpha}{1}$ \\
  $\eps\,(w\cdot\nabla\alpha_2)$ & $\tfrac12\eps^2t^2\,M_2\nrm{\nabla\alpha_2}{2}$ \\
  $\eps\theta\,(w\cdot\nabla)^2\alpha_2$ & $\tfrac1{12}\eps^3t^4\,M_4^2\nrm{\nabla^2\alpha_2}{2}$ \\
  $\eps^3\theta^3d\,(w\cdot\nabla\alpha_2)$ & $\tfrac1{20}\eps^4t^5\,d\,M_2\nrm{\nabla\alpha_2}{2}$ \\
  $\eps^2\theta\,(G\cdot\nabla\alpha_2)$ & $\tfrac12\eps^2t^2\,\nrm{|G||\nabla\alpha_2|}{1}$ \\
  $\eps^2\theta^2(w\cdot\nabla\alpha_2)g$ & $\tfrac1{12}\eps^3t^4\,M_2\nrm{|\nabla\alpha_2||g|}{2}$ \\
  $2\eps^2\theta\beta$ & $\eps^2t^2\,\nrm{\beta}{1}$ \\
  $\eps^2\theta^2(w\cdot\nabla\beta)$ & $\tfrac1{12}\eps^3t^4\,M_2\nrm{\nabla\beta}{2}$ \\
  $\eps^4\theta^4d\,\beta$ & $\tfrac15\eps^4t^5\,d\,\nrm{\beta}{1}$ \\
  $\eps^3\theta^3\beta g$ & $\tfrac14\eps^3t^4\,\nrm{\beta g}{1}$ \\
\end{tabular}
\end{center}
Collecting powers gives \eqref{eq:main} with
\begin{equation}
\begin{aligned}
  C_1 &= \nrm{\alpha}{1}, \\
  C_2 &= \tfrac12M_2\nrm{\nabla\alpha_2}{2}
        + \tfrac12\nrm{|G||\nabla\alpha_2|}{1} + \nrm{\beta}{1}, \\
  C_3 &= \tfrac14d\nrm{\alpha}{1} + \tfrac1{12}M_4^2\nrm{\nabla^2\alpha_2}{2}
        + \tfrac1{12}M_2\nrm{|\nabla\alpha_2||g|}{2}
        + \tfrac1{12}M_2\nrm{\nabla\beta}{2} + \tfrac14\nrm{\beta g}{1}, \\
  C_4 &= \tfrac1{20}d\,M_2\nrm{\nabla\alpha_2}{2} + \tfrac15d\nrm{\beta}{1}.
\end{aligned}
\label{eq:consts}
\end{equation}
With $t = s\eps^{-2/3}$: $\eps t = s\eps^{1/3}$,
$\eps^2t^2 = s^2\eps^{2/3}$, $\eps^3t^4 = s^4\eps^{1/3}$,
$\eps^4t^5 = s^5\eps^{2/3}$. \qed

\begin{remark}
Only the trivial bound $|w_\tau|\lesssim\eps\tau$ is used, not the diffusive
bound $\eps\tau^{1/2}$. The corrector $\alpha_2$ exists to remove the one term,
$i\eps\theta(w\cdot\nabla\alpha)U$, that the trivial bound cannot control at
$t\sim\eps^{-2/3}$. The term $\eps^3t^4$ is what limits the range to
$t\ll N^{3/8}$; a diffusive a priori bound would extend it.
\end{remark}

\section{The first iterate and the two thresholds}

\begin{proposition}\label{prop:first}
Let $Y^{(1)}_t = \eps\int_0^t(t-\sigma)\,G(0,b_\sigma)\,\dd\sigma$ be the
first Duhamel iterate of \eqref{eq:Yw}. Then $q\cdot Y^{(1)}_t$ is a centred
Gaussian variable with
\[
  \operatorname{Var}(q\cdot Y^{(1)}_t)
  = \eps^2\Bigl(\tfrac23\,d\,t^3 + O(t^2)\Bigr).
\]
In contrast, for a deterministic field with $|F|\le1$ the first iterate can
be as large as $\tfrac12\eps t^2$.
\end{proposition}

\begin{proof}
$g(0,b_\sigma)$ is a stationary Gaussian process with covariance
$c(u) = \frac12\sum_k|q\cdot f_k|^2\Re(e^{-\lambda_k|u|})$, and
$\int_0^\infty c = d$. By symmetry,
\[
  \operatorname{Var} = 2\eps^2\int_0^t c(u)\int_0^{t-u}a(a+u)\,\dd a\,\dd u
  = 2\eps^2\int_0^t c(u)\Bigl[\tfrac13(t-u)^3 + \tfrac12u(t-u)^2\Bigr]\dd u ,
\]
and the bracket equals $\tfrac13t^3 + O(t^2u + u^3)$ while $c$ decays
exponentially. The worst case is attained by a constant field.
\end{proof}

So the worst-case and mean-square sizes of the first iterate are $\eps t^2$
and $\eps t^{3/2}$, which reach order one at $t\sim N^{1/4}$ and $N^{1/3}$.
Theorem~\ref{thm:main} shows that at $N^{1/3}$ the full nonlinear orbit has
the same law as the first iterate to leading order, although the second
iterate is by then of the same size as the first: the perturbation series
does not converge there, but its sum is the Gaussian law of the first term.

\section{Numerical check}

A Monte Carlo evaluation of $\E e^{iq\cdot Y_t}$ was run on
16 November 2026 as a check on the algebra: $d=2$, five modes with
$f_k = -i\varphi_k k$, $q = (1,0.5)$, $q\cdot Dq = 0.9278$, $40000$ samples,
time step $0.02$, exact Ornstein--Uhlenbeck updates. It is a scratch run, not
part of the repository.
\begin{center}
\begin{tabular}{rrrrrr}
  $N$ & $s$ & $t$ & measured (real part) & predicted & MC s.e. \\ \hline
  $10^3$ & 1.0 & 10.0 & 0.7542 & 0.7340 & 0.0033 \\
  $10^3$ & 1.5 & 15.0 & 0.3741 & 0.3521 & 0.0046 \\
  $10^4$ & 1.0 & 21.5 & 0.7413 & 0.7340 & 0.0034 \\
  $10^4$ & 1.5 & 32.3 & 0.3636 & 0.3521 & 0.0047 \\
  $10^5$ & 1.0 & 46.4 & 0.7385 & 0.7340 & 0.0034 \\
  $10^5$ & 1.5 & 69.6 & 0.3598 & 0.3521 & 0.0047 \\
\end{tabular}
\end{center}
The discrepancy decreases with $N$ and is well inside the $O(N^{-1/6})$
bound; at $N = 10^5$ it is within about two standard errors. Imaginary parts
were within two standard errors of zero except at $N=10^3$, $s=1$.

\section{What separates the model from the particle system}
\label{sec:gap}

\begin{enumerate}
  \item \textbf{Self-consistency.} The field is produced by the other
    particles, whose orbits are themselves perturbed. One intermediate case
    is known: a tagged particle in a non-interacting equilibrium gas, in
    dimension $d\ge4$ (Bodineau and Le Bris, arXiv:2602.16440).
  \item \textbf{Back-reaction.} The tagged particle polarizes the medium,
    which produces a friction of order $1/N$. On $t\sim N^{1/3}$ it shifts
    the position by $t^2/N\sim N^{-1/3}$, so it should not change the
    statement, but this has to be proved.
  \item \textbf{No Markov structure.} The proof uses the generator $A$ of
    the mode amplitudes and solves Poisson equations for it. In the particle
    system the corresponding object is the linearized Vlasov flow, whose
    decay is Landau damping and holds only in weak norms.
  \item \textbf{Sharp resonances.} Here each mode has width $\gamma_k>0$. In
    the particle system the resonance is $\delta(k\cdot(v-v_*))$ integrated
    over a continuum of velocities; the constants $C_i$ above grow as
    $\gamma_k\to0$.
  \item \textbf{Non-Gaussian data.} Shot noise has cumulants of all orders.
    The proof uses only moments of the stationary law of $b$, so this is
    probably the mildest of the five.
  \item \textbf{From decorrelation to closure.} Even with the particle
    version of Theorem~\ref{thm:main}, steps 6 and 7 of the strategy (the
    bootstrap for the remainder, and iteration over layers) remain.
\end{enumerate}

\end{document}
